% \section{Related Work}
\subsection{Approaches for code localization}
Several prior methods have developed methods targeting code localization for downstream issue resolution in repositories. \citet{jimenez2024swebenchlanguagemodelsresolve} identify relevant source code files using BM25 retrieval~\citep{Robertson2009ThePR} treating the issue description as the query and the Python source code files as documents. SWE-Fixer~\citep{xie-etal-2025-swe} use a two-step coarse-to-fine localization approach where they first retrieve relevant source code files using BM25 retrieval and then prompt a fine-tuned 7B model with the skeleton of the retrieved files to predict the names of relevant files for issue resolution. On the other hand, Agentless~\citep{xia2024agentless} uses a complex, multi-step localization pipeline wherein it first retrieves suspicious code files from the repository using both embedding-based retrieval and by prompting an LLM with the high-level repository structure, followed by prompting an LLM to predict relevant functions and classes given the skeletons of the suspicious files. 

Many recent approaches have increasingly shifted towards developing models and systems for code localization using agentic scaffolds. Table~\ref{tab:granular_tool_behaviours} provides an overview of tools used by several prior agentic approaches and \methodname for code localization. LocAgent~\citep{chen-etal-2025-locagent} and RepoGraph~\citep{ouyang2025repographenhancingaisoftware} utilize a graph-based indexing approach wherein all dependencies in the code repository (for e.g. import, invoke, inherit, etc.) are captured in a code graph and also develop specialized tools allowing the agent to search and traverse the graph. Furthermore, \citet{chen-etal-2025-locagent} train a 7B model for their scaffold using rejection sampling fine-tuning on successful trajectories sampled from Claude-3.5-Sonnet. On the other hand, CoSIL~\citep{jiang2025issuelocalizationllmdriveniterative} first localizes relevant files using module-call graphs and then identifies relevant functions using function-call graphs. Note that these call graphs are constructed on-the-fly during inference as opposed to pre-indexing the code graph done by ~\citet{chen-etal-2025-locagent}. OrcaLoca~\citep{yu2025orcalocallmagentframework} performs code localization through efficient exploration of the code graph using priority-based action scheduling, action decomposition with relevance scoring, and distance-aware context pruning. Note that many of these scaffolds are not suitable for reinforcement learning as they often require expensive repository pre-processing (for e.g. creating a code graph) increasing the computational overhead of performing rollouts during reinforcement learning.

RepoSearcher~\citep{ma2025toolintegratedreinforcementlearningrepo} develop a light-weight scaffold with tools specialized for localization that allow searching for classes and functions in files, obtaining imports of a given file, etc. They use a two-stage approach to train open-source LLMs for their scaffold: (1) rejection sampling fine-tuning from a closed-source LLM (Claude3.7-Sonnet) to warm-up their model, (2) reinforcement learning on the fine-tuned checkpoint to further enhance performance. RepoNavigator~\citep{zhang2025one} use an even simpler scaffold comprising with just one tool - \texttt{jump} - which allows the agent to retrieve definitions of Python symbols in files. Although~\citet{zhang2025one} do not rely on rejection sampling fine-tuning from closed-source LLMs and directly train models using reinforcement learning, they still rely on selecting ``easy" training instances by discarding all training instances that were not successfully solved by the base Qwen2.5-7B model equipped with their scaffold atleast once among 16 sampled trajectories. 

\subsection{Training methods for software engineering agents}
Several prior approaches have trained LLM-based software engineering agents for various downstream tasks like code localization and issue resolution. Many methods rely on performing rejection sampling fine-tuning~\citep{yuan2023scalingrelationshiplearningmathematical}: sample trajectories (for training instances) from a stronger (often closed-source) model and train a smaller model on trajectories which successfully solve the task. Prior approaches that use this approach include LocAgent~\citep{chen-etal-2025-locagent} and RepoSearcher~\citep{ma2025toolintegratedreinforcementlearningrepo} for code localization, and SWE-Gym~\citep{pan2025trainingsoftwareengineeringagents}, SWE-Smith~\citep{yang2025swesmithscalingdatasoftware}, and R2E-Gym~\citep{jain2025r2egymproceduralenvironmentshybrid} for issue resolution. Recently, various attempts have been made to train models directly with reinforcement learning without relying on proprietary, closed-source LLMs for rejection sampling fine-tuning. While RepoNavigator~\citep{zhang2025one}, SWE-Grep~\citep{cognition2025swegrep}, and SID-1~\citep{sid2025preview} train LLM agents with RL for code localization, DeepSWE~\citep{luo2025deepswe} and SkyRL-v0~\citep{cao2025skyrl} leverage RL to train LLM agents for issue resolution.

% \subsection{Multi-agent and modular systems}
% While most prior approaches primarily focus on 

% Modular architectures decompose complex software engineering tasks into specialized components. MASAI~\citep{arora2024masaimodulararchitecturesoftwareengineering} employs multiple subagents for issue reproduction, localization, fixing, and ranking. In industry practice, specialized subagents for retrieval and context gathering have also become common~\citep{cognition2025swegrep,morph2025warpgrep,cherny2025claudecode,heule2025semsearch}. Such modularity motivates training dedicated localization components that can be integrated with downstream editing agents.

% \subsection{Positioning}
% Our work targets multi-level localization (file/class/function) with a lightweight tool scaffold and trains a dedicated localization subagent directly with RL. Compared to structure-heavy approaches (e.g., graph-guided localization) and general repository-level agents, we emphasize a simple, reproducible recipe with multi-level F1 rewards and stability techniques, and release the full training and evaluation pipeline.
